\section{Refraction is conformal geometry, not a new law}\label{sec:refraction}

The corpus states a refraction law $\rho_1\sin^2\theta_1=\rho_2\sin^2\theta_2$ across a
``rigidity interface''. The law is true, and it is Snell's law for a conformal metric; saying so
precisely is the whole content of this section, and it changes the epistemic status of the claim
from a transport law to a consequence of a modelling choice.

\begin{theorem}[Snell from a conformal factor]\label{thm:snell}
Let $g$ be a Riemannian metric on a surface, $R>0$ a function that is $\C^1$ on each side of a
$\C^1$ curve $\shape$ and jumps across it, and $\hat g=R\,g$. Let $\gamma$ be a geodesic of
$\hat g$ crossing $\shape$ transversally at $p$. Then
\begin{equation}\label{eq:snell}
   n_1\sin\theta_1=n_2\sin\theta_2,\qquad\text{equivalently}\qquad
   R_1\sin^2\theta_1=R_2\sin^2\theta_2,
\end{equation}
where $n:=\sqrt R$ and $\theta$ is the angle between $\gamma$ and the normal to $\shape$
measured in the metric $g$.
\T\ (first variation of the $\hat g$-length, with the matching condition at $\shape$).
\end{theorem}

\begin{proof}
A geodesic of $\hat g$ is a critical point of $\hat g$-length against variations supported away
from the endpoints. Taking a variation localised near $p$ and letting its support shrink gives
continuity across $\shape$ of the tangential component of the $\hat g$-unit tangent $u$: for every
vector $\tau$ tangent to $\shape$, $\ip{u}{\tau}_{\hat g}$ is continuous. Let $\tau$ be the
$g$-unit tangent and write $\abs{u}_g=R^{-1/2}$ because $\ip{u}{u}_{\hat g}=R\ip{u}{u}_g=1$.
Then $\ip{u}{\tau}_{\hat g}=R\ip{u}{\tau}_g=R\abs{u}_g\sin\theta=\sqrt R\,\sin\theta$, so
$\sqrt{R_1}\sin\theta_1=\sqrt{R_2}\sin\theta_2$; squaring gives \eqref{eq:snell}.
\end{proof}

\begin{remark}[What is a theorem and what is a hypothesis]\label{rem:snell-grade}
Theorem~\ref{thm:snell} is \T. It is not a statement about any physical transport: it is a
statement about geodesics of a conformal metric. To apply it to a transport one must assert that
admissible paths extremise $\int R\sin^2\theta\,\df s$, or equivalently that the effective metric
is conformal to a reference one with conformal factor $R$. That assertion is a modelling choice,
and it is \Hh\ wherever the corpus makes it. With $\rho$ in the role of $R$ one recovers the
corpus's law and the index $n=\sqrt\rho$; with $\rho=2/(3-\sqrt{9-4\det})$ of
Remark~\ref{rem:rigidity} one gets a specific, falsifiable prediction for the shape sphere,
namely an index that diverges like $1/A$ as collinearity is approached, so that transport across
a near-collinear interface is totally reflected to leading order.
\end{remark}

\begin{remark}[Why the qualification matters]
A ``refraction law'' sounds like a discovered regularity of a system. \eqref{eq:snell} is instead
what one gets for free once one has decided that the medium is conformal with factor $R$ --- the
same content as Fermat's principle, re-expressed. The empirical question is never whether the
law holds, but whether the medium is conformal and with which $R$. That is the question a gate can
answer, and it is the form in which the corpus's interface-refraction protocol should be stated.
\end{remark}
